\section{Related Work}
\label{sec:related}

%\akshita{add the new related works; or maybe we sprinkle them in earlier in intro}

\paragraph{Scaling laws for language modeling.}
\citet{Kaplan2020ScalingLF} and \citet{Hoffmann2022TrainingCL} were among the first to postulate the functional form of LM losses as a power function of model parameters and data size.
The Chinchilla equation, in particular, has become the basis of many subsequent scaling law works \citep{Muennighoff2023ScalingDL,Gadre2024LanguageMS,Zhang2024MAPNeoHC}.
This line of work focuses on predicting the LM loss (rather than downstream performance) on a held-out set with similar distribution as the pretraining data, but not on downstream tasks which is more important and challenging.

\tightparagraph{Scaling laws for downstream tasks.}
Scaling laws for downstream tasks have been explored in \citet{Gadre2024LanguageMS}. Instead of directly predicting accuracy for individual tasks, they compute the average top-1 error over 17 LLM-foundry \citep{mosaicml_llm_foundry} evaluation tasks as a function of cross entropy loss on C4 \citep{Raffel2019ExploringTL}. \citet{Dubey2024TheL3} uses a two-step prediction to first map the the training compute to the negative log-likelihood of the correct answer for a single task in an evaluation benchmark, and then relate the log-likelihood to the task accuracy. Unlike our work, which uses a fixed ladder of small models, they rely on first finding compute-optimal models. % Their target model (405B model trained to 15T tokens) is also in the compute-optimal regime, whereas, we aim to predict models trained beyond compute-optimality. \ljc{This sentence feels a bit out-of-place to me, it breaks the flow of presenting existing approaches to downstream tasks.}
\citet{Chen2024ScalingLF} also employs a two-stage approach for predicting downstream performance, but uses the pre-training loss instead of a task-specific loss as the intermediate step.
%\citet{Hu2023PredictingEA} predicts performance on generative tasks by introducing an intermediate ``PassUntil'' metric; the number of repeated decoding required to land on a correct answer.
\citet{Isik2024ScalingLF} studies scaling laws in the finetuning setup, for machine translation tasks. \citet{polo2025slothscalinglawsllm} leverages models of different families to predict downstream performance, but this leads to less accurate predictions for guiding pretraining development for a specific family of models.
